\documentclass[conference]{IEEEtran}
\IEEEoverridecommandlockouts
\usepackage{cite}
\usepackage{amsmath,amssymb,amsfonts}
\usepackage{algorithmic}
\usepackage{graphicx}
\usepackage{textcomp}
\usepackage{xcolor}
\usepackage{hyperref}

\def\BibTeX{{\rm B\kern-.05em{\sc i\kern-.025em b}\kern-.08em
    T\kern-.1667em\lower.7ex\hbox{E}\kern-.125emX}}

\begin{document}

\title{Evaluasi Peningkatan Citra Berbasis HVI-CIDNet untuk Deteksi Objek Low-Light pada ExDark dengan YOLOv11}

\author{\IEEEauthorblockN{Muhammad Iqbal Rasyid}
\IEEEauthorblockA{\textit{Dept. of Informatics} \\
\textit{Telkom University}\\
Bandung, Indonesia \\
otachiking@student.telkomuniversity.ac.id}
}

\maketitle

\begin{abstract}
Kualitas citra pada kondisi minim cahaya (low-light) menjadi tantangan besar untuk tugas deteksi objek dalam visi komputer. Sinyal lemah yang ditangkap oleh sensor kamera mengakibatkan tingkat derau (noise) yang tinggi dan kontras yang rendah, sehingga menyulitkan model deteksi untuk mengenali fitur objek dengan baik. Untuk mengatasi tantangan tersebut, penelitian ini mengevaluasi integrasi metode peningkatan citra (image enhancement) berbasis Color and Intensity Decoupling Network (CIDNet) yang beroperasi pada ruang warna HVI, dengan arsitektur deteksi objek mutakhir, YOLOv11. Evaluasi dilakukan menggunakan dataset Exclusively Dark (ExDark) yang mencakup 12 kelas objek dalam berbagai kondisi pencahayaan minim nyata. Hasil dari penelitian ini diharapkan dapat memberikan wawasan empiris mengenai seberapa besar pengaruh pra-pemrosesan HVI-CIDNet terhadap peningkatan akurasi metrik mean Average Precision (mAP) dari model YOLOv11.
\end{abstract}

\begin{IEEEkeywords}
Low-Light Image Enhancement, Object Detection, HVI-CIDNet, YOLOv11, ExDark
\end{IEEEkeywords}

\section{Introduction}
Kinerja algoritma deteksi objek telah mengalami kemajuan pesat dalam dekade terakhir, terutama dengan hadirnya keluarga model \textit{You Only Look Once} (YOLO) \cite{terven2023yoloreview}. Meskipun model mutakhir ini menunjukkan performa luar biasa pada citra dengan pencahayaan normal, akurasi deteksi seringkali terdegradasi ketika diterapkan pada citra dengan kondisi minim cahaya (\textit{low-light}). Tantangan utama pada lingkungan \textit{low-light} meliputi visibilitas yang buruk, kontras yang sangat rendah, dan hilangnya detail penting akibat tingginya tingkat \textit{noise} bawaan sensor kamera \cite{loh2019exdark}. Untuk mengatasi masalah ini, pendekatan \textit{Low-Light Image Enhancement} (LLIE) konvensional umumnya diposisikan sebagai langkah pra-pemrosesan (\textit{preprocessing}) wajib sebelum citra dimasukkan ke dalam model deteksi \cite{liu2021benchmark}.

Baru-baru ini, berbagai metode LLIE tingkat lanjut seperti \textit{Color and Intensity Decoupling Network} (CIDNet) yang terintegrasi dengan ruang warna HVI (Hue, Value, Illumination), RetinexFormer, dan LYT-Net telah menunjukkan kemampuan luar biasa dalam memulihkan kecerahan citra sekaligus menekan \textit{noise} \cite{yan2025hvi}, \cite{cai2023retinexformer}, \cite{brateanu2025lytnet}. Namun, terdapat perdebatan akademis yang krusial mengenai efektivitas algoritma ini pada tugas visi tingkat tinggi (\textit{high-level machine vision}). Peningkatan kualitas visual yang memanjakan persepsi mata manusia terbukti tidak selalu berbanding lurus dengan peningkatan kinerja ekstraksi fitur oleh detektor mesin \cite{wu2024llieeffect}. Algoritma LLIE rentan mengintroduksi \textit{noise} frekuensi tinggi serta mengubah struktur semantik fitur asli citra selama proses restorasi, yang justru berpotensi bersifat kontraproduktif terhadap akurasi deteksi \cite{wang2023yolov5lowlight}.

Di sisi lain, arsitektur YOLOv11 hadir sebagai detektor objek sekali-jalan (\textit{one-stage detector}) terbaru yang menawarkan kemampuan ekstraksi fitur yang sangat tangguh. Melalui integrasi blok C3k2, \textit{Spatial Pyramid Pooling-Fast} (SPPF), dan mekanisme atensi spasial C2PSA, YOLOv11 mampu mengekstraksi representasi secara langsung dari data mentah (\textit{raw data}) yang minim cahaya dengan sangat efektif \cite{khanam2024yolov11}. Kecanggihan struktural ini memunculkan hipotesis empiris bahwa pemberian input citra yang telah dimanipulasi oleh algoritma LLIE justru berisiko mendisorientasi persepsi jaringan akibat hilangnya gradien asli serta penambahan artefak buatan dari proses restorasi warna.

Berdasarkan fenomena tersebut, penelitian ini menyajikan sebuah studi empiris untuk mengevaluasi secara komprehensif pengaruh pra-pemrosesan metode LLIE SOTA terhadap kinerja deteksi objek YOLOv11. Pengujian difokuskan pada perbandingan metrik kinerja ekstraksi fitur dan akurasi (mAP) antara citra mentah sebagai \textit{baseline} melawan citra yang diproses menggunakan metode HVI-CIDNet, RetinexFormer, dan LYT-Net. Dataset \textit{Exclusively Dark} (ExDARK) dipilih sebagai basis pengujian karena representasinya yang kuat terhadap variasi skenario pencahayaan rendah di dunia nyata \cite{loh2019exdark}. Evaluasi ini dilakukan menggunakan skema inferensi langsung tanpa penyesuaian ulang parameter detektor untuk mengisolasi variabel eksperimen dan menjaga objektivitas. Selain memvalidasi degradasi performa secara kuantitatif, penelitian ini juga berkontribusi secara analitis dengan menggunakan metode interpretabilitas \textit{Mean Activation Map} \cite{metrics2026formulas}. Pendekatan visualisasi tersebut diterapkan untuk membuktikan secara konkret bagaimana pergeseran fokus komputasi dan distorsi peta fitur terjadi pada struktur internal detektor akibat paparan artefak dari metode peningkatan citra.

\section{Related Work}
Evaluasi kinerja deteksi objek pada lingkungan minim cahaya melibatkan irisan dari dua domain penelitian utama, yakni algoritma perbaikan citra dan arsitektur visi mesin tingkat tinggi. Bagian ini menguraikan perkembangan terkini pada kedua domain tersebut serta celah penelitian empiris yang menghubungkan keduanya.

\subsection{Low-Light Image Enhancement (LLIE)} Tujuan utama dari algoritma LLIE adalah memulihkan visibilitas, kontras, dan warna dari citra yang terdegradasi akibat minimnya iluminasi \cite{liu2021benchmark}. Secara historis, metode ini berevolusi dari pendekatan tradisional seperti manipulasi histogram menjadi model berbasis Deep Learning yang kompleks. Pendekatan berbasis Vision Transformer seperti RetinexFormer \cite{cai2023retinexformer} dan arsitektur ringan LYT-Net \cite{brateanu2025lytnet} telah membuktikan kemampuan superior dalam mengekstraksi fitur pencahayaan global.

Lebih lanjut, metode mutakhir mulai mengeksplorasi manipulasi ruang warna untuk menghindari distorsi yang umum terjadi pada ruang sRGB dan HSV \cite{yan2025hvi}. Pendekatan Color and Intensity Decoupling Network (CIDNet) yang beroperasi pada ruang Horizontal/Vertical-Intensity (HVI) terbukti secara efektif mampu memisahkan pemrosesan intensitas cahaya dan warna secara independen \cite{yan2025hvi}. Modul Lighten Cross-Attention (LCA) pada arsitektur ini mencegah amplifikasi noise silang antara saluran warna dan kecerahan \cite{darmawan2025vitra}. Meskipun algoritma LLIE mutakhir ini menghasilkan skor kuantitatif yang sangat tinggi pada metrik perseptual manusia seperti Peak Signal-to-Noise Ratio (PSNR) dan Structural Similarity (SSIM), optimasi perseptual ini tidak dirancang secara spesifik untuk mempertahankan gradien fitur yang krusial bagi visi mesin.

\subsection{Evolusi Arsitektur You Only Look Once (YOLO)} Keluarga algoritma YOLO telah mendominasi tugas deteksi objek real-time berkat arsitekturnya yang efisien dalam memprediksi bounding box dan probabilitas kelas secara serentak \cite{terven2023yoloreview}, \cite{sapkota2025yoloreview}. Pada generasi terdahulu seperti YOLOv3 dan YOLOv5, arsitektur ekstraksi fitur memiliki keterbatasan dalam menangani citra dengan rasio signal-to-noise yang sangat rendah \cite{wang2023yolov5lowlight}. Oleh karena itu, penerapan metode LLIE sebagai pra-pemrosesan terbukti secara empiris mampu mendongkrak performa (Mean Average Precision) pada model-model generasi lama tersebut \cite{shovo2024lowlight}, \cite{peng2024yolov5llie}.
Namun, iterasi terbaru yakni YOLOv11 memperkenalkan pembaruan struktural yang masif pada bagian backbone dan neck. YOLOv11 mengintegrasikan blok Cross Stage Partial with kernel size 2 (C3k2) untuk optimasi ekstraksi fitur secara komputasional, modul Spatial Pyramid Pooling-Fast (SPPF) untuk ketahanan skala, serta Convolutional block with Parallel Spatial Attention (C2PSA) yang secara dramatis memperkuat mekanisme atensi spasial jaringan terhadap fitur objek \cite{khanam2024yolov11}. Kecanggihan arsitektural ini membuat model mampu membedakan fitur objek dari noise latar belakang secara mandiri, bahkan pada data citra mentah (raw) yang sangat gelap.

\subsection{Dampak LLIE terhadap Visi Mesin Tingkat Tinggi} Terdapat diskoneksi fundamental antara optimasi citra untuk persepsi visual manusia dan optimasi citra untuk analisis mesin. Studi empiris komprehensif yang dilakukan oleh Wu et al. \cite{wu2024llieeffect} mengungkapkan bahwa penerapan algoritma LLIE dapat berdampak inkonsisten pada tugas klasifikasi dan deteksi, bahkan sering kali bersifat merusak (harmful).

Proses pemulihan iluminasi dan warna pada LLIE rentan menyuntikkan noise frekuensi tinggi, artefak buatan, serta pergeseran semantik yang secara permanen mengubah distribusi fitur gradien asli citra \cite{wang2023yolov5lowlight}, \cite{gong2025multiscale}. Bagi detektor mutakhir seperti YOLOv11 yang memiliki arsitektur ekstraksi fitur sangat sensitif dan tangguh, intervensi dari artefak LLIE justru dapat mendisorientasi fokus atensi komputasinya. Alih-alih meningkatkan performa, pra-pemrosesan ini berpotensi menurunkan akurasi deteksi lokalisasi objek. Oleh karena itu, penelitian yang mengevaluasi secara kritis dampak negatif dari metode LLIE mutakhir (seperti HVI-CIDNet) terhadap detektor modern di dataset nyata seperti ExDARK \cite{loh2019exdark} menjadi sangat relevan. Evaluasi ini membutuhkan pembuktian tidak hanya melalui metrik performa akhir, tetapi juga melalui visualisasi peta aktivasi fitur (Mean Activation Map) guna mengobservasi pergeseran lokalisasi fitur di dalam struktur jaringan detektor \cite{metrics2026formulas}.

% ============================================== %
% [PLACEHOLDER_GAMBAR_2: Diagram alir atau ilustrasi komparatif konseptual yang menunjukkan diskoneksi antara output LLIE (Human Vision) dengan ekstraksi fitur oleh YOLOv11 (Machine Vision)]
% BERIKUT ADALAH LAYOUT UNTUK BAB-BAB SELANJUTNYA %
% ============================================= %

\section{Methodology}
Penelitian ini menggunakan kerangka kerja eksperimental komparatif untuk mengevaluasi dampak prapemrosesan algoritma Low-Light Image Enhancement (LLIE) terhadap detektor objek YOLOv11. Kerangka kerja ini dibagi menjadi empat tahapan utama, yakni persiapan dataset, inferensi prapemrosesan LLIE, pelatihan model deteksi, dan metode interpretabilitas fitur visual.

\subsection{Dataset and Preprocessing} Evaluasi empiris dilakukan menggunakan dataset publik Exclusively Dark (ExDARK) yang dirancang khusus untuk skenario pencahayaan rendah \cite{loh2019exdark}. Dataset ini mencakup 7.363 citra dunia nyata yang terbagi ke dalam 12 kategori objek (seperti \textit{Bicycle}, \textit{Car}, \textit{People}, dan lainnya) serta 10 kondisi variasi iluminasi. Untuk memastikan validitas pelatihan model, dataset direstrukturisasi menjadi tiga subset utama secara acak, yakni 3.000 citra untuk pelatihan (\textit{training}), 1.800 citra untuk validasi (\textit{validation}), dan 2.563 citra untuk pengujian akhir (\textit{testing}) \cite{proposal2026sempro}.
Seluruh anotasi bounding box dikonversi dari format asli ke dalam format standar YOLO dengan normalisasi koordinat spasial \cite{padilla2021metrics}, \cite{proposal2026sempro}. Untuk memenuhi standar arsitektur YOLOv11 tanpa merusak proporsi fitur geometris objek, seluruh citra disesuaikan dimensinya menjadi 640x640 piksel menggunakan metode \textit{letterbox resizing}, yang mempertahankan rasio aspek asli dengan menambahkan padding pada sisi citra \cite{proposal2026sempro}.

\subsection{Low-Light Image Enhancement Pipeline} Untuk mengevaluasi dampak manipulasi intensitas cahaya terhadap visi mesin, penelitian ini membandingkan skenario citra mentah (\textit{baseline}) dengan tiga skenario prapemrosesan metode LLIE \textit{state-of-the-art} (SOTA), yakni HVI-CIDNet \cite{yan2025hvi}, RetinexFormer \cite{cai2023retinexformer}, dan LYT-Net \cite{brateanu2025lytnet}.
Ketiga metode tersebut diaplikasikan menggunakan strategi inferensi langsung (\textit{direct inference}) tanpa melakukan penyesuaian ulang parameter (\textit{fine-tuning}) pada dataset LOLv1 \cite{proposal2026sempro}. Model LLIE dijalankan menggunakan bobot pra-latih (\textit{pretrained weights}) asli untuk menguji kemampuan generalisasi algoritma secara objektif. Prapemrosesan HVI-CIDNet secara spesifik memisahkan pemrosesan intensitas dan warna melalui transformasi ruang warna Horizontal/Vertical-Intensity (HVI) sebelum mengembalikan citra hasil restorasi ke dalam format sRGB \cite{yan2025hvi}, \cite{proposal2026sempro}. Citra keluaran dari ketiga metode inilah yang kemudian dikonfigurasi sebagai input dataset sekunder untuk pelatihan model detektor.

\subsection{Object Detection Architecture and Training} Arsitektur detektor yang digunakan difokuskan pada varian YOLOv11n (Nano) \cite{khanam2024yolov11}. Pemilihan varian dengan kompleksitas terendah ini bertujuan untuk mengisolasi variabel eksperimen, karena model berkapasitas kecil memiliki sensitivitas yang lebih tinggi terhadap kualitas gradien input dibandingkan model berkapasitas besar \cite{proposal2026sempro}. Hal ini memastikan bahwa fluktuasi akurasi deteksi murni disebabkan oleh kualitas fitur citra, bukan oleh ketahanan (\textit{robustness}) internal model.
Pelatihan model YOLOv11n pada masing-masing skenario dilakukan secara konsisten menggunakan metrik dari laporan evaluasi \cite{report2026internal}. Pelatihan dijalankan selama 50 \textit{epochs} dengan ukuran batch 16. Optimasi jaringan menggunakan algoritma AdamW dengan \textit{learning rate} awal sebesar 0.005 dan \textit{weight decay} 0.001. Untuk meningkatkan generalisasi detektor, strategi augmentasi data diaktifkan dengan konfigurasi Mosaic (1.0), Mixup (0.2), Erasing (0.2), serta manipulasi parameter ruang warna HSV \cite{report2026internal}. Ambang batas validasi kepercayaan (\textit{confidence threshold}) diatur pada 0.001 dengan metrik Intersection over Union (IoU) sebesar 0.7 \cite{report2026internal}, \cite{metrics2026formulas}.

\subsection{Evaluation Metrics and Model Interpretability} Kinerja dari setiap skenario diukur menggunakan standar evaluasi kuantitatif COCO, yang mencakup presisi, recall, serta \textit{Mean Average Precision} pada ambang batas IoU 0.5 (mAP@0.5) dan rentang komprehensif 0.5 hingga 0.95 (mAP@0.5:0.95) \cite{report2026internal}, \cite{metrics2026formulas}. Kompleksitas sistem divalidasi menggunakan kalkulasi \textit{Giga Floating-Point Operations} (GFLOPs) dan \textit{Frames Per Second} (FPS) \cite{metrics2026formulas}.

Untuk menginvestigasi secara kualitatif perdebatan empiris mengenai degradasi fitur akibat introduksi artefak LLIE, penelitian ini mengimplementasikan metode interpretabilitas model berbasis \textit{PyTorch Forward Hook} tanpa perhitungan gradien balik \cite{metrics2026formulas}. Visualisasi ekstraksi fitur diamati menggunakan dua pendekatan:

\textbf{\textit{Mean Activation Map}}: Diterapkan pada \textit{Layer} 0 untuk mengevaluasi respons jaringan terhadap fitur tingkat rendah (seperti tepi dan tekstur) serta pada \textit{Layer} N untuk mengevaluasi fitur semantik. Peta ini dihitung dengan merata-ratakan seluruh aktivasi saluran jaringan \cite{metrics2026formulas}.

\textbf{EigenCAM (\textit{Class-Agnostic Saliency})}: Diterapkan pada \textit{Layer} N menggunakan proyeksi dekomposisi \textit{Principal Component} pertama (PC1) dari matriks aktivasi spasial. Karena bersifat agnostik terhadap kelas, EigenCAM secara objektif memvisualisasikan disorientasi atensi jaringan utama akibat \textit{noise} bawaan dari algoritma LLIE tanpa perlu melakukan adaptasi gradien per kelas objek \cite{metrics2026formulas}.

% ============================================== %
% [PLACEHOLDER_GAMBAR_3: Diagram skematik yang mengilustrasikan keseluruhan alur metodologi, mulai dari input ExDARK, pemisahan cabang (Raw vs LLIE), hingga ke tahap deteksi YOLOv11 dan ekstraksi fitur (EigenCAM)]
% ============================================== %

\section{Results and Discussion}
Bagian ini menyajikan hasil evaluasi empiris mengenai dampak algoritma Low-Light Image Enhancement (LLIE) terhadap ekstraksi fitur dan kinerja detektor YOLOv11. Analisis dibagi menjadi evaluasi kualitas citra spasial, dinamika pelatihan, kinerja deteksi kuantitatif, serta interpretasi model secara kualitatif.

\subsection{Visual Quality and Spatial Feature Degradation} Evaluasi awal berfokus pada analisis dampak prapemrosesan LLIE (Skenario S2 hingga S4) terhadap struktur spasial citra dibandingkan dengan citra mentah atau baseline (Skenario S1). Meskipun metode LLIE berhasil secara signifikan menaikkan nilai rata-rata iluminasi (hingga 114.59 pada HVI-CIDNet) dan memperbaiki metrik perseptual buta (No-Reference Metrics) seperti NIQE dan BRISQUE \cite{report2026internal}, optimasi visual ini harus dibayar dengan degradasi fitur struktural yang parah.

Berdasarkan kalkulasi metrik spasial, prapemrosesan LLIE secara konsisten mengintroduksi noise baru dan merusak konsistensi batas objek (edge). Skenario S1 memiliki tingkat noise terendah ($\sigma$ = 4.515). Sebaliknya, algoritma HVI-CIDNet (S2), RetinexFormer (S3), dan LYT-Net (S4) masing-masing meningkatkan tingkat noise menjadi 7.427, 6.833, dan 7.136 \cite{report2026internal}. Kerusakan struktural ini paling jelas terlihat pada metrik Edge Preservation Index (EPI). Nilai EPI yang ideal adalah 1.0 (seperti pada S1). Namun, setelah prapemrosesan LLIE, nilai EPI anjlok drastis ke kisaran 0.2004 (S2) hingga 0.2696 (S4) \cite{report2026internal}, \cite{metrics2026formulas}. Penurunan drastis pada EPI ini mengindikasikan bahwa algoritma LLIE telah menggeser, menghilangkan, atau menciptakan artefak tepi palsu yang sangat destruktif bagi proses lokalisasi bounding box pada visi mesin.

% [PLACEHOLDER_GAMBAR_1: Perbandingan Antara Gambar Original vs Gambar Enhanced (S2-S4). Catatan untuk Anda: Gunakan gambar yang memperlihatkan dengan jelas bagaimana S2-S4 menjadi lebih terang tetapi memunculkan noise atau artefak aneh di latar belakang]

\subsection{Training Dynamics and Convergence} Dinamika pelatihan YOLOv11n pada keempat skenario dievaluasi untuk mengamati stabilitas konvergensi model selama 50 epochs. Seluruh metrik kerugian (Box Loss, Classification Loss, dan Distribution Focal Loss atau DFL) dimonitor secara ketat.

Pelatihan menggunakan dataset baseline (S1) menunjukkan kurva penurunan kerugian (loss) yang paling stabil dan mulus pada fase validasi. Sebaliknya, pelatihan menggunakan citra dari skenario S2 hingga S4 memperlihatkan fluktuasi loss yang lebih tinggi, terutama pada Box Loss dan DFL Loss. Fluktuasi ini mengonfirmasi hipotesis bahwa pergeseran nilai piksel spasial dan artefak buatan dari algoritma LLIE membuat jaringan YOLOv11 kesulitan untuk mempelajari batas kotak pembatas yang konsisten.

% [PLACEHOLDER_GAMBAR_2: Ultralytic Training Summary (gabungan train/box_loss, train/cls_loss, val/box_loss, dll) DAN Training Curves (Precision, Recall, mAP@0.5). Catatan: Bisa digabung dalam satu grid multi-panel]

\subsection{Quantitative Object Detection Performance} Evaluasi kinerja deteksi akhir membuktikan bahwa optimasi kecerahan untuk persepsi manusia tidak sejalan dengan kebutuhan ekstraksi fitur oleh YOLOv11. Skenario S1 (Raw) secara konsisten mengungguli seluruh varian citra yang telah ditingkatkan. Hal ini sejalan dengan temuan Wu et al. \cite{wu2024llieeffect} bahwa model dengan ekstraksi fitur superior justru mengalami disorientasi saat diberikan citra manipulasi.

Hasil evaluasi kuantitatif menunjukkan bahwa distribusi presisi dan recall pada skenario S1 lebih tinggi dibandingkan metode SOTA seperti HVI-CIDNet \cite{yan2025hvi} dan RetinexFormer \cite{cai2023retinexformer}. Selain itu, kurva F1-Confidence dan Precision-Confidence menunjukkan bahwa model baseline (S1) mampu mempertahankan tingkat kepercayaan (kestabilan prediksi) yang lebih tinggi pada rentang klasifikasi 12 kelas objek ExDARK. Matriks kebingungan (Confusion Matrix) juga memperlihatkan bahwa tingkat False Positive pada skenario S2 hingga S4 meningkat secara signifikan. Peningkatan kesalahan ini utamanya disebabkan oleh model yang mendeteksi tumpukan noise sisa algoritma restorasi warna sebagai fitur objek palsu.

% [PLACEHOLDER_GAMBAR_3: mAP Progression (perbandingan mAP0.5 dan mAP0.5:0.95 lintas skenario)]
% [PLACEHOLDER_GAMBAR_4: F1-Confidence Curve dan Precision-Confidence Curve (versi 12 Class + All)]
% [PLACEHOLDER_GAMBAR_5: Confusion Matrix (Versi Count dan Normalized untuk menonjolkan perbedaan False Positives)]

\subsection{Feature Degradation and Model Interpretability} Untuk menguji secara kualitatif penyebab penurunan akurasi pada citra yang ditingkatkan (enhanced), penelitian ini menerapkan metode interpretabilitas PyTorch Forward Hook melalui visualisasi Mean Activation Map dan EigenCAM pada lapisan akhir (Layer N) YOLOv11 \cite{metrics2026formulas}.

Visualisasi prediksi bounding box (Prediksi versus Groundtruth) memperlihatkan bahwa model pada skenario S2 hingga S4 sering kali gagal melakukan pelingkupan (bounding) objek secara presisi, meskipun objek terlihat sangat terang bagi mata manusia. Analisis peta fitur mengonfirmasi fenomena tersebut. Pada citra mentah (S1), peta Mean Activation dan EigenCAM menunjukkan bahwa konsentrasi atensi jaringan YOLOv11 terpusat secara akurat pada fitur semantik objek target. Sebaliknya, pada citra hasil HVI-CIDNet (S2) dan varian LLIE lainnya, atensi komputasi model pecah dan terdistraksi (disoriented) ke area latar belakang. Jaringan detektor merespons artefak perbaikan warna dan noise frekuensi tinggi sebagai "salien fitur", yang pada akhirnya merusak pemahaman konteks spasial secara keseluruhan.

% [PLACEHOLDER_GAMBAR_6: Detection Results: Prediction vs Groundtruth. Catatan: Pilih contoh gambar di mana S1 berhasil mendeteksi objek dengan pas, sementara S2-S4 kotaknya meleset atau ada false positive]

% [PLACEHOLDER_GAMBAR_7: Model Interpretability: Original, First Layer, Mean Activation, dan Eigen-CAM. Catatan: Ini adalah "Golden Image" Anda yang membuktikan atensi model hancur (tersebar ke background) akibat efek LLIE]

\section{Conclusion}
Penelitian ini telah melakukan evaluasi empiris secara komprehensif mengenai dampak penerapan algoritma \textit{Low-Light Image Enhancement} (LLIE) \textit{state-of-the-art}, yakni HVI-CIDNet \cite{yan2025hvi}, RetinexFormer \cite{cai2023retinexformer}, dan LYT-Net \cite{brateanu2025lytnet}, terhadap kinerja detektor objek YOLOv11 pada dataset nyata ExDARK \cite{loh2019exdark}. Berlawanan dengan asumsi konvensional bahwa peningkatan visibilitas visual selalu menguntungkan sistem deteksi, hasil eksperimen membuktikan bahwa arsitektur YOLOv11 justru mencapai tingkat akurasi rata-rata (mAP) yang lebih superior dan stabil ketika memproses citra \textit{raw low-light} secara langsung tanpa pra-pemrosesan apa pun.
Penurunan performa pada skenario citra yang ditingkatkan (\textit{enhanced}) terbukti secara kuantitatif disebabkan oleh introduksi \textit{noise} frekuensi tinggi serta kerusakan pelestarian tepi spasial, yang tercermin dari penurunan drastis pada metrik \textit{Edge Preservation Index} (EPI). Arsitektur YOLOv11, yang telah dilengkapi dengan mekanisme ekstraksi fitur mutakhir seperti blok C3k2 dan atensi spasial C2PSA \cite{khanam2024yolov11}, terbukti memiliki kemampuan persepsi \textit{raw} yang kuat namun sangat rentan terhadap artefak buatan dari algoritma restorasi warna. Analisis kualitatif menggunakan metode interpretabilitas \textit{Mean Activation Map} dan EigenCAM \cite{metrics2026formulas} secara konkret mengonfirmasi fenomena tersebut. Pra-pemrosesan LLIE mendisorientasi persepsi jaringan sehingga atensi komputasional yang seharusnya terpusat pada fitur semantik objek justru terpecah dan merespons \textit{noise} latar belakang sebagai fitur palsu.
Temuan ini menyoroti adanya diskoneksi fundamental antara algoritma perbaikan citra yang dioptimalkan untuk persepsi mata manusia dengan kebutuhan pelestarian gradien murni untuk visi mesin \cite{wu2024llieeffect}. Oleh karena itu, arah penelitian di masa depan sangat direkomendasikan untuk beralih dari optimasi metrik perseptual (seperti PSNR atau SSIM) menuju pengembangan metode pemulihan yang berorientasi pada pelestarian semantik (\textit{task-driven enhancement}). Lebih lanjut, eksplorasi optimasi gabungan secara \textit{end-to-end} atau integrasi modul atensi sadar iluminasi ke dalam struktur internal detektor berpotensi memberikan solusi yang lebih tangguh untuk sistem deteksi objek pada lingkungan minim cahaya ekstrem.

\bibliographystyle{IEEEtran}
\bibliography{referensi} % Memanggil file referensi.bib

\end{document}